\chapter{Модулярный угол кривизны спина три и соответствие переходов}

Предыдущая глава нашла единственную семёрку спина три внутри
$\operatorname{Hom}(V_1,V_2)$ и показала, что её квадрат на триплетном
углу равен тетраэдрической кривизне. Буквальный самосопряжённый блок
$3\leftrightarrow5$ был закрыт двумерным коядром. Здесь проверяется,
может ли прежняя модулярная архитектура соответствий канонически выбрать
нужный односторонний угол, не превращая весь прямоугольный блок в спектр
физических состояний.

\section{Антикруговой журнал}

Гейт не повторяет обычный спектральный квадрат. Он уже закрыт, поскольку
видит оба диагональных угла. Не повторяется и произвольная угловая
компрессия: ранний запрет выбора удобного угла требует, чтобы проектор
был задан состоянием, градуировкой или правилом суперселекции до
вычисления вакуума.

Проект уже содержит подходящий механизм. В модулярном соответствии
положительная частота оператора $\operatorname{ad}_h$ выбирает
ориентированный дифференциал, а условная нормировка спектрального
проектора вершины восстанавливает её нормированный след. Это было
получено в гейтах модулярного соответствия и углового чтения состояния.

В общей литературе правое гильбертово $C^*$-модульное скалярное
произведение является алгеброзначным и естественно односторонним.
Кривизна дифференцируемых гильбертовых модулей относительно спектральной
тройки построена в \path{arXiv:1911.05008}; общая теория бимодульных
связностей и кривизн --- в \path{arXiv:q-alg/9512004}. Эти источники
обосновывают тип конструкции, но не выбирают проектную семёрку и её
нормировку.

\section{Связывающий контейнер и модулярная высота}

Пусть
\begin{equation}
 E=\operatorname{Hom}(V_1,V_2),
 \qquad
 \mathcal L(E)=
 \begin{pmatrix}
  \operatorname{End}(V_1)&E^*\\
  E&\operatorname{End}(V_2)
 \end{pmatrix}
 \simeq M_8(\mathbb R).
 \label{eq:v6-spin-three-linking-container}
\end{equation}
Обозначим проекторы двух углов через $p_1$ и $p_2$, их ранги равны трём
и пяти. Высота
\begin{equation}
 h=-p_1+p_2
 \label{eq:v6-spin-three-modular-height}
\end{equation}
коммутирует с диагональным действием одной семейной группы $SO(3)$.

Для стрелки $Z:V_1\to V_2$ введём
\begin{equation}
 d=
 \begin{pmatrix}0&0\\ Z&0\end{pmatrix}.
\end{equation}
Тогда
\begin{equation}
 [h,d]=2d,
 \qquad
 D=d+d^*.
 \label{eq:v6-spin-three-positive-modular-frequency}
\end{equation}
Остатки обоих тождеств равны нулю. Следовательно, $p_1$ является
исходным проектором положительно-частотной стрелки и выбирается до задания
тетраэдрического вакуума. Он не подбирается по знаку полученного
потенциала.

\begin{proposition}[Модулярное происхождение исходного угла]
Высота~\eqref{eq:v6-spin-three-modular-height} и положительная частота
однозначно различают источник $V_1$ и цель $V_2$ ориентированной стрелки.
Поэтому компрессия на $p_1$ является частью данных соответствия, а не
послевакуумным выбором удобного угла.
\end{proposition}

\section{Односторонняя кривизна}

Градуированный коммутатор равен
\begin{equation}
 [d,d^*]=
 \begin{pmatrix}
  -Z^*Z&0\\0&ZZ^*
 \end{pmatrix}.
\end{equation}
Поэтому исходная кривизна определяется формулой
\begin{equation}
 \mathcal F_1
 =-p_1[d,d^*]p_1-\frac{v_T^2}{3}p_1
 =Z^*Z-\frac{v_T^2}{3}I_3.
 \label{eq:v6-modular-source-tetrahedral-curvature}
\end{equation}
После ограничения $Z=Z_T$ на каноническую семёрку предыдущей главы
получаем точное равенство
\begin{equation}
 \mathcal F_1=\mu_T(T).
\end{equation}
На тетраэдрическом вакууме остаток равен нулю.

Проектор Казимира $P_{(3)}$ коммутирует со всеми тремя генераторами
физического $SO(3)$ с максимальным остатком
$2.3\cdot10^{-15}$. Инфинитезимальная ковариантность квадрата имеет
нулевой остаток:
\begin{equation}
 \delta(Z^*Z)=[\omega,Z^*Z].
\end{equation}
Таким образом, семимерное пространство является корректным
ассоциированным подрасслоением одной семейной связности.

\section{Единственный след и фиксированный угловой вес}

Связывающий контейнер имеет один нормированный след
\begin{equation}
 \tau_8(A)=\frac18\Tr_8A,
 \qquad
 \tau_8(p_1)=\frac38.
\end{equation}
Условная нормировка исходного угла даёт
\begin{equation}
 \frac{\tau_8(p_1\mathcal F_1^2p_1)}{\tau_8(p_1)}
 =\frac13\Tr_3\mathcal F_1^2.
 \label{eq:v6-spin-three-conditioned-source-trace}
\end{equation}
Ненормированное чтение полного контейнера отличается ровно фиксированным
ранговым множителем $3/8$. На тридцати двух случайных полях максимальные
остатки условного тождества и ранговой формулы равны соответственно нулю
и $4.5\cdot10^{-16}$.

После усиления физическим пакетом $H_{15}$ связывающий носитель имеет
размерность
\begin{equation}
 (3+5)\cdot15=120,
\end{equation}
а исходный проектор --- ранг сорок пять. Поэтому его доля снова равна
$45/120=3/8$, и условный след точно восстанавливает уже доказанный след
на $H_{45}$. Максимальный остаток кривизны равен
$8.9\cdot10^{-16}$, кинетической нормы --- $1.8\cdot10^{-15}$.

\begin{theorem}[Односторонний родительский квадрат]
Модулярно ориентированное соответствие $V_1\to V_2$ вместе с единственным
следом связывающего контейнера канонически даёт действие
\begin{equation}
 S_T=\frac13\Tr_3
 \left(Z_T^*Z_T-\frac{v_T^2}{3}I_3\right)^2
\end{equation}
без нового относительного веса. После усиления $H_{15}$ нормировка
совпадает с физическим следом $H_{45}$.
\end{theorem}

\section{Что стало с двумя нулевыми направлениями}

Два нуля полного блока $D=d+d^*$ не исчезают алгебраически. Меняется их
роль. В текущей ветви $D$ не является новым конечным оператором физических
фермионов и не входит в спектральный определитель. Он служит
самосопряжённым завершением ориентированной стрелки соответствия.
Действие строится из кривизны исходного угла, поэтому коядро $ZZ^*$ не
считается набором новых состояний.

Это не переоткрывает закрытую буквальную дираковскую ветвь. Если
$D$ всё же объявить физическим конечным оператором, два нуля и их
тридцатикратное усиление возвращаются немедленно.

\begin{nogocriterion}[Разделение оператора и соответствия]
Нельзя одновременно использовать одностороннюю модульную норму и считать
весь восьмимерный самосопряжённый блок новым физическим спектром. Первый
вариант является бозонным соответствием; второй остаётся закрыт
коядровым препятствием.
\end{nogocriterion}

\section{Ограничение полной контейнерной алгебры}

Семёрка $\operatorname{im}P_{(3)}$ сохраняется физическим $SO(3)$, но
не замкнута относительно произвольного правого умножения на $M_3$.
Случайный тест даёт минимальный выход из семёрки $0.9344\ldots$.
Поэтому её нельзя объявлять гильбертовым подмодулем над всей $M_3$.

Это согласуется со старым запретом калибровки связывающего контейнера.
Полное пространство $E$ задаёт операторную среду, а физическое поле
$Z_T$ является ассоциированным $SO(3)$-подрасслоением. Полная $M_3$ не
становится новой координатной или калибровочной алгеброй.

\section{Вердикт}

\begin{protocol}
\begin{enumerate}
 \item модулярная высота: \textbf{получена};
 \item исходный угол выбран до вакуума: \textbf{да};
 \item односторонняя кривизна: \textbf{$\mu_T$ точно};
 \item ковариантность относительно одной $SO(3)$: \textbf{получена};
 \item условный след из единственного $\tau_8$: \textbf{получен};
 \item ранговый вес исходного угла: \textbf{$3/8$};
 \item усиление до $H_{120}$ и ограничение на $H_{45}$:
 \textbf{совместимы};
 \item новый относительный коэффициент: \textbf{нет};
 \item семёрка как полный правый $M_3$-модуль: \textbf{нет};
 \item полная $M_3$ как физическая калибровочная алгебра: \textbf{нет};
 \item буквальная конечная дираковская ветвь: \textbf{остаётся закрытой};
 \item бозонная ветвь соответствия на уровне родительской нормы:
 \textbf{проходит};
 \item рождение материи: \textbf{не закрыто};
 \item следующий гейт: \textbf{совместный вакуум $Q+T$ и полной
 семейной связности}.
\end{enumerate}
\end{protocol}

Машинный сертификат сохранён в
\path{s2t_v6_bosonic_defect_spin_three_corner_curvature_parent_gate_results.json}.